\chapter{Multivariate Series and Cointegration}\label{ch:qm:multivariate-series-and-cointegration}

The two-year, five-year and ten-year Treasury yields each look like random walks: over fifty years of daily data, none of the three rejects
a \omterm{def:qm:linear-time-series:unitroot}{unit root} (Dickey--Fuller statistics of $-1.30$, $-1.33$ and $-1.31$). A curve fly that is long the five-year and short the wings does not
wander off in the same way: the yields share their trends, and some combination of the three is stationary. The question a relative-value
desk asks is which weights make it so, and how fast the resulting fly comes back. The textbook fly weights the wings one-two-one; the data prefer
to put 0.41 on the two-year and 0.61 on the ten-year for each unit of the five-year, and that fly reverts with a \omterm{def:qm:stochastic-differential-equations:ou}{half-life} of 43 trading days
against 63. This chapter is about several series at once: \omterm{def:qm:multivariate-series-and-cointegration:var}{vector autoregressions} and what their shocks do, \omterm{def:qm:multivariate-series-and-cointegration:coint}{cointegration} and the error-correction
form, the Engle--Granger and \omterm{def:qm:multivariate-series-and-cointegration:johansen}{Johansen tests} that find it, and causality in the predictive sense.

\section{Vector autoregressions}

\begin{definition}[Vector autoregression]\label{def:qm:multivariate-series-and-cointegration:var}
A \emph{vector autoregression}\index{vector autoregression} of order $p$, VAR($p$), models a $k$-vector as $Y_t = c + \sum_{i=1}^pA_iY_{t-i} + \varepsilon_t$ with
$\varepsilon_t$ \omterm{def:qm:linear-time-series:stationary}{white noise} of covariance $\Sigma$ (Sims, 1980). Each equation is estimated by least squares on the same regressors.
\end{definition}

\begin{proposition}[Stationarity of a VAR]\label{prop:qm:multivariate-series-and-cointegration:stable}
The VAR is weakly stationary if and only if every eigenvalue of its companion matrix $\begin{psmallmatrix}A_1 & \cdots & A_{p-1} & A_p\\ I & & & 0\\ & \ddots & & \vdots\\
& & I & 0\end{psmallmatrix}$ has modulus below one; equivalently, $\det(I - \sum_iA_iz^i) \ne 0$ for $|z| \le 1$.
\end{proposition}

\begin{proof}
Stack $(Y_t, \dots, Y_{t-p+1})$ into a VAR(1) with the companion matrix $C$. Its moving-average coefficients are $C^h$, which are summable exactly when the
spectral radius of $C$ is below one; the eigenvalues of $C$ are the reciprocals of the roots of the determinant.
\end{proof}

\begin{definition}[Impulse response function, forecast-error variance decomposition]\label{def:qm:multivariate-series-and-cointegration:irf}
The \emph{impulse response function}\index{impulse response function} gives the effect on $Y_{t+h}$ of a one-standard-deviation shock at $t$; with correlated
shocks, they are first orthogonalised by the Cholesky factor of $\Sigma$, which makes the result depend on the ordering of the variables. The
\emph{forecast-error variance decomposition}\index{forecast-error variance decomposition} gives the share of each variable's $h$-step forecast-error variance
attributable to each orthogonalised shock.
\end{definition}

On the daily yield changes (in basis points), AIC keeps improving up to the ten lags it is allowed, and the VAR(10) is stable. The residuals have standard
deviations of 7.9, 7.6 and 7.1 basis points and are strongly correlated (0.90 between the two- and five-year, 0.82 between the two- and ten-year, 0.94 between
the five- and ten-year), so the ordering matters. With the two-year first, a one-standard-deviation two-year shock of 7.9 basis points moves the two-, five-
and ten-year yields by a cumulative 10.1, 8.3 and 6.7 basis points over ten days (\cref{fig:qm:multivariate-series-and-cointegration:irf}), and it accounts
for 67\% of the ten-year's ten-day forecast-error variance (the five-year's shock for 21\%, the ten-year's own for 12\%). The two-year leads, in this ordering,
by construction; whether it leads in time is a separate question.

\begin{omfigure}
\begin{tikzpicture}
\begin{axis}[width=10cm,height=5cm,xlabel={days after the shock},ylabel={cumulative response (bp)},
  tick label style={font=\small},label style={font=\small},xmin=0,xmax=10,ymin=0,ymax=11,grid=major,grid style={gray!25},
  legend columns=3,legend style={font=\footnotesize,at={(0.5,-0.3)},anchor=north,draw=none,fill=none,
  /tikz/every even column/.append style={column sep=8pt}},table/col sep=comma]
\addplot[omDef,very thick,mark=*,mark size=1.3pt] table[x=day,y=y2]{figdata/methods/20-multivariate-series-and-cointegration/irf.csv};
\addplot[black,very thick,mark=square*,mark size=1.3pt] table[x=day,y=y5]{figdata/methods/20-multivariate-series-and-cointegration/irf.csv};
\addplot[omThm,very thick,mark=triangle*,mark size=1.6pt] table[x=day,y=y10]{figdata/methods/20-multivariate-series-and-cointegration/irf.csv};
\legend{2-year,5-year,10-year}
\end{axis}
\end{tikzpicture}

\omcaption{Cumulative orthogonalised responses of the 2-, 5- and 10-year yields to a one-standard-deviation (7.9 bp) shock to the 2-year, from a VAR(10) of
daily changes, 1976--2026, ordered 2-, 5-, 10-year. Data: FRED series DGS2, DGS5 and DGS10 (Board of Governors of the Federal Reserve System,
H.15).}\label{fig:qm:multivariate-series-and-cointegration:irf}
\end{omfigure}

\section{Causality in the predictive sense}

\begin{definition}[Granger causality]\label{def:qm:multivariate-series-and-cointegration:granger}
A series $X$ \emph{Granger-causes}\index{Granger causality} $Y$ if the past of $X$ improves the prediction of $Y$ beyond the past of $Y$ and of the other
variables (Granger, 1969); in a VAR, if the coefficients of the lags of $X$ in $Y$'s equation are not all zero, tested by an $F$-statistic.
\end{definition}

It is predictive precedence, not causation: a common driver seen earlier in one series than in the other produces it, and so does a timing difference
in how the data are recorded. On the daily yield changes, every pair \omterm{def:qm:multivariate-series-and-cointegration:granger}{Granger-causes} the other at 5\%, with $F(10, 12\,532)$ between 1.88 (two-year to
ten-year, $p = 0.04$) and 7.16 (two-year to five-year, $p < 10^{-6}$); the ten-year to the two-year gives 2.83 ($p = 0.002$). With 12\,574 days, tiny lead-lag
effects are significant; none is large enough to trade on its own.

\section{Cointegration and the error-correction form}

\begin{definition}[Cointegration, cointegrating vector, cointegration rank]\label{def:qm:multivariate-series-and-cointegration:coint}
The components of a vector of unit-root processes are cointegrated (\emph{cointegration}\index{cointegration}) if some linear combination $\beta^\top Y_t$ is
stationary; $\beta$ is a \emph{cointegrating vector}\index{cointegrating vector}. The number of linearly independent cointegrating vectors is the
\emph{cointegration rank}\index{cointegration rank} $r$, and $k - r$ is the number of common stochastic trends.
\end{definition}

\begin{definition}[Vector error-correction model]\label{def:qm:multivariate-series-and-cointegration:vecm}
A \emph{vector error-correction model}\index{vector error-correction model} writes a VAR in differences plus a levels term, $\Delta Y_t = \alpha\beta^\top Y_{t-1} +
\sum_{i=1}^{p-1}\Gamma_i\Delta Y_{t-i} + \mu + \varepsilon_t$, with $\alpha$ and $\beta$ of dimension $k \times r$: $\beta^\top Y_{t-1}$ is the deviation from equilibrium
and $\alpha$ the speed at which each variable corrects it.
\end{definition}

By the Granger representation theorem (Engle and Granger, 1987), a cointegrated system has an error-correction representation and conversely: a VAR in levels
whose matrix $\Pi = \sum_iA_i - I$ has reduced rank $r$ factors as $\alpha\beta^\top$. A VAR in differences alone throws the levels information away; a VAR in
levels keeps it but hides it. For the yields, $k = 3$; if their curve is driven by level and slope trends with a stationary curvature, $r = 1$ and the
\omterm{def:qm:multivariate-series-and-cointegration:coint}{cointegrating vector} is a fly.

\begin{definition}[Engle--Granger test]\label{def:qm:multivariate-series-and-cointegration:eg}
The \emph{Engle--Granger test}\index{Engle--Granger test} regresses one variable on the others in levels and applies a \omterm{def:qm:linear-time-series:df}{Dickey--Fuller test} to the residuals,
with critical values that account for the estimated coefficients.
\end{definition}

The least-squares residual is the most stationary-looking combination of the sample, so it looks stationary more often than a fixed one: its Dickey--Fuller
statistic is shifted left, and the critical values depend on the number of variables. For three variables MacKinnon's (2010) 5\% value is $-3.74$, against
$-2.86$ for one series; in the chapter's simulation of 4\,000 triples of independent random walks, using the one-series value would find \omterm{def:qm:multivariate-series-and-cointegration:coint}{cointegration} 27\% of the
time (\cref{fig:qm:multivariate-series-and-cointegration:egdf}). Regressing the five-year yield on the other two gives 0.40 times the two-year plus 0.62 times
the ten-year, and a residual statistic of $-10.04$: \omterm{def:qm:multivariate-series-and-cointegration:coint}{cointegration}, overwhelmingly.

\begin{omfigure}
\begin{tikzpicture}
\begin{axis}[width=10cm,height=5cm,xlabel={Dickey--Fuller $t$-statistic},ylabel={density under the null},
  tick label style={font=\small},label style={font=\small},xmin=-6,xmax=2,ymin=0,ymax=0.6,grid=major,grid style={gray!25},
  legend columns=2,legend style={font=\footnotesize,at={(0.5,-0.3)},anchor=north,draw=none,fill=none,
  /tikz/every even column/.append style={column sep=8pt}},table/col sep=comma]
\addplot[omThm,thick,const plot mark mid,fill=omThm!15] table[x=tau,y=df]{figdata/methods/20-multivariate-series-and-cointegration/egdf.csv};
\addplot[omDef,very thick,const plot mark mid] table[x=tau,y=eg]{figdata/methods/20-multivariate-series-and-cointegration/egdf.csv};
\addplot[omThm,thin,dashed,sharp plot] coordinates {(-2.92,0) (-2.92,0.6)};
\addplot[omDef,thin,dashed,sharp plot] coordinates {(-3.75,0) (-3.75,0.6)};
\legend{one random walk,Engle--Granger residual (3 series)}
\end{axis}
\end{tikzpicture}

\omcaption{Null distributions of the Dickey--Fuller statistic for one random walk and for the least-squares residual of one random walk regressed on two
others, from 4\,000 simulated samples of 1\,000 days; the dashed lines are their 5\% points, $-2.92$ and $-3.75$. Data: the chapter's tutorial,
seeded.}\label{fig:qm:multivariate-series-and-cointegration:egdf}
\end{omfigure}

\section{Rank tests}

\begin{definition}[Johansen test]\label{def:qm:multivariate-series-and-cointegration:johansen}
The \emph{Johansen test}\index{Johansen test} (Johansen, 1988, 1991) estimates the VECM by maximum likelihood under the restriction $\operatorname{rank}\Pi = r$:
after removing the short-run dynamics from $\Delta Y_t$ and $Y_{t-1}$, the squared canonical correlations $\hat\lambda_1 \ge \dots \ge \hat\lambda_k$ between the two
residual sets give the trace statistic $-T\sum_{i>r}\ln(1 - \hat\lambda_i)$ for rank at most $r$, and the maximum-eigenvalue statistic $-T\ln(1 - \hat\lambda_{r+1})$;
the corresponding eigenvectors estimate $\beta$.
\end{definition}

The null distributions are functionals of \omterm{def:qm:brownian-motion:bm}{Brownian motion} that depend on the number of common trends and on the deterministic terms. With the constant
restricted to the cointegrating space (yields have no trend), the chapter simulates them from 20\,000 systems of random walks of 2\,000 days; the 95\% points of the
trace statistic are 35.23, 20.36 and 9.23 for three, two and one common trend.

On the three yields, 1976--2026, with nine lagged differences (the VAR(10) of AIC), the trace statistics are 59.77 for rank zero, 14.93 for rank at most one
and 2.34 for rank at most two: rank one, a single \omterm{def:qm:multivariate-series-and-cointegration:coint}{cointegrating vector}. Normalised on the five-year yield, it is $0.41 \times$ two-year $- 1 \times$ five-year $+ 0.61
\times$ ten-year: a fly whose wings sum to 1.02, close to level-neutral, but tilted toward the ten-year, so that it also neutralises most of the slope. As an AR(1),
it reverts with coefficient 0.9841, a \omterm{def:qm:stochastic-differential-equations:ou}{half-life} of 43 trading days; the one-two-one fly (scaled to the same five-year weight) has coefficient 0.9891 and a \omterm{def:qm:stochastic-differential-equations:ou}{half-life}
of 63 days, and a daily volatility of 2.15 basis points against 2.06 (\cref{fig:qm:multivariate-series-and-cointegration:flies}).

\begin{omfigure}
\begin{tikzpicture}
\begin{groupplot}[group style={group size=1 by 2,vertical sep=0.5cm,xticklabels at=edge bottom},width=12cm,height=4.2cm,
  tick label style={font=\small},label style={font=\small},xmin=1976,xmax=2026.8,grid=major,grid style={gray!25},
  xticklabel style={/pgf/number format/set thousands separator={}},table/col sep=comma]
\nextgroupplot[ylabel={yield (bp)},ymin=0,ymax=1700,legend columns=3,
  legend style={font=\footnotesize,at={(0.98,0.96)},anchor=north east,fill=white,draw=none}]
\addplot[omDef,thick] table[x=year,y=y2]{figdata/methods/20-multivariate-series-and-cointegration/yields.csv};
\addplot[black,thick] table[x=year,y=y5]{figdata/methods/20-multivariate-series-and-cointegration/yields.csv};
\addplot[omThm,thick] table[x=year,y=y10]{figdata/methods/20-multivariate-series-and-cointegration/yields.csv};
\legend{2-year,5-year,10-year}
\nextgroupplot[xlabel={year},ylabel={fly (bp, demeaned)},ymin=-80,ymax=80,legend columns=2,
  legend style={font=\footnotesize,at={(0.98,0.96)},anchor=north east,fill=white,draw=none}]
\addplot[omThm!70,thin] table[x=year,y=fly121]{figdata/methods/20-multivariate-series-and-cointegration/yields.csv};
\addplot[omDef,thin] table[x=year,y=flyj]{figdata/methods/20-multivariate-series-and-cointegration/yields.csv};
\legend{one-two-one,Johansen}
\end{groupplot}
\end{tikzpicture}

\omcaption{Top: the 2-, 5- and 10-year Treasury yields, 1976--2026, each indistinguishable from a random walk. Bottom: two flies, long the 5-year and short the
wings, with the 5-year weight one: the one-two-one weights (0.5, 0.5) and the Johansen weights (0.41, 0.61), whose half-lives are 63 and 43 trading days. Data:
FRED series DGS2, DGS5 and DGS10 (H.15).}\label{fig:qm:multivariate-series-and-cointegration:flies}
\end{omfigure}

The weights are an estimate, and not a stable one. On five-year rolling windows, the trace test finds no \omterm{def:qm:multivariate-series-and-cointegration:coint}{cointegration} at 5\% in 20 of the 46 windows, and in the
windows where it finds one vector, the two-year weight ranges from 0.27 to 0.63 and the ten-year weight from 0.40 to 0.97
(\cref{fig:qm:multivariate-series-and-cointegration:rolling}). Fifty years pin the fly down; five years do not, and a desk that trades the fly lives in the five years.

\begin{omfigure}
\begin{tikzpicture}
\begin{axis}[width=11cm,height=5cm,xlabel={last year of the five-year window},ylabel={fly weight},
  tick label style={font=\small},label style={font=\small},xmin=1980,xmax=2027,ymin=0,ymax=1.2,grid=major,grid style={gray!25},
  xticklabel style={/pgf/number format/set thousands separator={}},
  legend columns=2,legend style={font=\footnotesize,at={(0.5,-0.3)},anchor=north,draw=none,fill=none,
  /tikz/every even column/.append style={column sep=8pt}},table/col sep=comma]
\addplot[only marks,mark=*,mark size=1.8pt,omDef] table[x=end,y=w2]{figdata/methods/20-multivariate-series-and-cointegration/rolling.csv};
\addplot[only marks,mark=square*,mark size=1.8pt,omThm] table[x=end,y=w10]{figdata/methods/20-multivariate-series-and-cointegration/rolling.csv};
\addplot[omDef,dashed,thin,sharp plot] coordinates {(1980,0.411) (2027,0.411)};
\addplot[omThm,dashed,thin,sharp plot] coordinates {(1980,0.607) (2027,0.607)};
\legend{2-year weight,10-year weight}
\end{axis}
\end{tikzpicture}

\omcaption{Johansen fly weights (per unit of the 5-year) estimated on rolling five-year windows, with the full-sample values 0.41 and 0.61 (dashed). Four windows
whose weights exceed 1.5 in absolute value are not shown; in 20 of the 46 windows the trace test does not find \omterm{def:qm:multivariate-series-and-cointegration:coint}{cointegration} at 5\%. Data: FRED series DGS2,
DGS5 and DGS10 (H.15).}\label{fig:qm:multivariate-series-and-cointegration:rolling}
\end{omfigure}

\section{Tutorial: the fly that mean-reverts}\label{tut:qm:multivariate-series-and-cointegration:fly}

\begin{tutorial}
\textbf{Goal.} Find the cointegrating fly of the 2-, 5- and 10-year yields and compare it with the one-two-one fly.
\textbf{End state:} \cref{fig:qm:multivariate-series-and-cointegration:flies,fig:qm:multivariate-series-and-cointegration:rolling}, rank one, the weights (0.41,
0.61) and the half-lives 43 and 63 days.
\begin{tutsteps}
\item \textbf{Johansen's reduced-rank regression}, with the constant restricted to the cointegrating space.
\omcode{code/firm/coint/firm_coint.py}{165}{198}{Johansen's canonical correlations, trace and maximum-eigenvalue statistics.}\label{lst:qm:multivariate-series-and-cointegration:johansen}
\item \textbf{Engle--Granger} with MacKinnon's response surfaces.
\omcode{code/firm/coint/firm_coint.py}{146}{162}{The Engle--Granger residual test.}\label{lst:qm:multivariate-series-and-cointegration:eg}
\item \textbf{Run} \texttt{cointegration()}, \texttt{var\_analysis()}, \texttt{rolling\_weights()} and \texttt{eg\_vs\_df()} in \texttt{qm\_coint.py}, then
\texttt{fig\_coint.py}.
\end{tutsteps}
\textbf{What to change next.} Weight the fly by duration (DV01) rather than by yield and redo the test on the P\&L; estimate the VECM's $\alpha$ and see which yield
does the correcting.
\end{tutorial}

\section{Build: the cointegration module}\label{bld:qm:multivariate-series-and-cointegration:coint}

\begin{build}
\bfield{Purpose} Relative-value baskets in the miniature firm (flies, spreads, index-versus-constituents) get their weights and their tests here.
\bfield{Interface} \texttt{var\_fit(Y, p)}, \texttt{var\_stable(A)}, \texttt{irf(A, Sigma, horizon)}, \texttt{fevd(A, Sigma, horizon)}; \texttt{granger\_test(Y,
p, cause, effect)}, \texttt{f\_sf}; \texttt{engle\_granger(y, X, lags)}; \texttt{johansen(Y, lags)}; \texttt{johansen\_crit\_sim(m, T, reps, seed)}.
\bfield{Rules} Critical values come from response surfaces or from this module's documented simulation, never from the Dickey--Fuller table; weights are reported
with the rank and the sample; the ordering of variables is stated with every orthogonalised response.
\bfield{Acceptance tests} \texttt{code/firm/coint/tests/}: a VAR(1) recovered, stability, the first impulse responses, decompositions summing to one; Granger tests with
power and correct size; the $F$ tail against the $\chi^2$ limit; Engle--Granger on a cointegrated pair and on independent walks; Johansen finds rank one and its
direction; the stored critical values re-simulate.
\bfield{Stretch} The VECM's $\alpha$ and $\Gamma$ with \omterm{def:qm:estimation:estimator}{standard errors}; weak exogeneity tests; the fly traded with a Kalman-filtered weight (chapter~19).
\end{build}

\begin{omsources}
\item C.~A.~Sims, ``Macroeconomics and reality'', \emph{Econometrica} 48, 1980.
\item C.~W.~J.~Granger, ``Investigating causal relations by econometric models and cross-spectral methods'', \emph{Econometrica} 37, 1969.
\item R.~F.~Engle and C.~W.~J.~Granger, ``Co-integration and error correction: representation, estimation, and testing'', \emph{Econometrica} 55, 1987.
\item S.~Johansen, ``Statistical analysis of cointegration vectors'', \emph{Journal of Economic Dynamics and Control} 12, 1988; ``Estimation and hypothesis
testing of cointegration vectors in Gaussian vector autoregressive models'', \emph{Econometrica} 59, 1991.
\item J.~G.~MacKinnon, ``Critical values for cointegration tests'', Queen's Economics Department Working Paper 1227, 2010.
\item Board of Governors of the Federal Reserve System, H.15 Selected Interest Rates, via FRED (series DGS2, DGS5 and DGS10), accessed 24 September 2026.
\end{omsources}

\section{Exercises}

\begin{exercise}[$\star$]\label{exo:qm:multivariate-series-and-cointegration:1}
Is the VAR(1) with $A = \begin{psmallmatrix}0.5 & 0.4\\ 0.4 & 0.5\end{psmallmatrix}$ stationary? And with $\begin{psmallmatrix}0.6 & 0.4\\ 0.4 & 0.6\end{psmallmatrix}$?
\end{exercise}

\begin{exercise}[$\star$]\label{exo:qm:multivariate-series-and-cointegration:2}
Two yields $x$ and $y$ are random walks with $y - 0.9x$ stationary. What are $k$, $r$ and the number of common trends?
\end{exercise}

\begin{exercise}[$\star$]\label{exo:qm:multivariate-series-and-cointegration:3}
Compute MacKinnon's 5\% Engle--Granger critical value for three variables and $T = 500$ from $(-3.74066, -8.5631, -10.852, 27.982)$.
\end{exercise}

\begin{exercise}[$\star\star$]\label{exo:qm:multivariate-series-and-cointegration:4}
Write the VECM of $y_t = 0.9x_{t-1} + 0.5(y_{t-1} - 0.9x_{t-1}) + e_t$ with $x$ a random walk, and identify $\alpha$ and $\beta$.
\end{exercise}

\begin{exercise}[$\star\star$]\label{exo:qm:multivariate-series-and-cointegration:5}
Two residuals have correlation 0.9 and unit variances. With the first variable ordered first, what fraction of the second variable's one-step forecast variance is
attributed to the first shock?
\end{exercise}

\begin{exercise}[$\star\star$]\label{exo:qm:multivariate-series-and-cointegration:6}
The Johansen eigenvalues of a three-variable system with $T = 1\,000$ are 0.05, 0.01 and 0.002. Compute the trace statistics and the rank at 5\%.
\end{exercise}

\begin{exercise}[$\star\star\star$]\label{exo:qm:multivariate-series-and-cointegration:7}
\emph{Coding.} Estimate the Johansen fly weights on rolling five-year windows of the yields, and count the windows in which the trace test finds no \omterm{def:qm:multivariate-series-and-cointegration:coint}{cointegration}.
\end{exercise}

\begin{exercise}[$\star\star\star$]\label{exo:qm:multivariate-series-and-cointegration:8}
\emph{Find the flaw.} ``We regressed the 5-year yield on the 2-year and 10-year, got a residual whose Dickey--Fuller statistic is $-3.1$, below the 5\% critical value
of $-2.86$, and started trading the residual as a stationary fly.''
\end{exercise}

\section{Problem: The Fly That Mean-Reverts}

\begin{problem}[{Weekend problem --- the cointegrating 2s5s10s fly}]\label{pb:qm:multivariate-series-and-cointegration:1}
The data are the daily 2-, 5- and 10-year constant-maturity Treasury yields (FRED DGS2, DGS5, DGS10), 1 June 1976 to 22 September 2026, 12\,574 days, in basis points.

\medskip\noindent\textbf{Part I --- Each series.}
\begin{enumerate}
\item What do \omterm{def:qm:linear-time-series:df}{Dickey--Fuller tests} say about each yield?
\item What VAR order does AIC choose for the daily changes, and is the VAR stable?
\item What are the residual standard deviations and correlations?
\item What does a one-standard-deviation 2-year shock do over ten days?
\item Which yields Granger-cause which, and what does it mean?
\end{enumerate}

\medskip\noindent\textbf{Part II --- \omterm{def:qm:multivariate-series-and-cointegration:coint}{Cointegration}.}
\begin{enumerate}[resume]
\item What are the Johansen trace statistics and the rank at 5\%?
\item What are the fly weights per unit of the 5-year, and what do they neutralise?
\item What does Engle--Granger give, against which critical value?
\item What are the half-lives of the Johansen fly and of the one-two-one fly?
\item What are the flies' daily volatilities?
\end{enumerate}

\medskip\noindent\textbf{Part III --- Robustness.}
\begin{enumerate}[resume]
\item How often does a five-year window find \omterm{def:qm:multivariate-series-and-cointegration:coint}{cointegration}?
\item How much do the weights move across windows?
\item Why is the Engle--Granger critical value lower than Dickey--Fuller's?
\item How would you correct the \omterm{def:qm:stochastic-differential-equations:ou}{half-life} for the small-sample bias of chapter~17?
\item What would a desk risk by trading the full-sample weights?
\end{enumerate}

\medskip\noindent\textbf{Part IV --- Judgement.}
\begin{enumerate}[resume]
\item Which fly would you trade, with which weights?
\item How would you size it for a \omterm{def:qm:stochastic-differential-equations:ou}{half-life} of 43 days?
\item Why should the weights be re-estimated, and how often?
\item State the \emph{named result}: the Johansen fly weights and \omterm{def:qm:stochastic-differential-equations:ou}{half-life} against the one-two-one fly's.
\item In one sentence: what does \omterm{def:qm:multivariate-series-and-cointegration:coint}{cointegration} promise a relative-value trader, and what not?
\end{enumerate}
\end{problem}

\section{Interview questions}

\begin{interviewq}[$\star$ \iqroles{researcher, trader}]\label{iq:qm:multivariate-series-and-cointegration:1}
What is the difference between correlation and \omterm{def:qm:multivariate-series-and-cointegration:coint}{cointegration}?
\end{interviewq}

\begin{interviewq}[$\star\star$ \iqroles{researcher, trader}]\label{iq:qm:multivariate-series-and-cointegration:2}
How would you choose the weights of a curve fly to trade \omterm{def:qm:stochastic-differential-equations:ou}{mean reversion}?
\end{interviewq}

\begin{interviewq}[$\star\star$ \iqroles{researcher}]\label{iq:qm:multivariate-series-and-cointegration:3}
Why can't you use the Dickey--Fuller critical values on the residual of a cointegrating regression?
\end{interviewq}

\begin{interviewq}[$\star\star$ \iqroles{researcher, mle}]\label{iq:qm:multivariate-series-and-cointegration:4}
What does \omterm{def:qm:multivariate-series-and-cointegration:granger}{Granger causality} test, and what does it not?
\end{interviewq}

\begin{interviewq}[$\star\star$ \iqroles{researcher}]\label{iq:qm:multivariate-series-and-cointegration:5}
Why does the ordering matter in a VAR's impulse responses?
\end{interviewq}

\begin{interviewq}[$\star\star\star$ \iqroles{researcher}]\label{iq:qm:multivariate-series-and-cointegration:6}
Explain the \omterm{def:qm:multivariate-series-and-cointegration:johansen}{Johansen test} in terms of canonical correlations.
\end{interviewq}
